\section{A stationary comparison with a labelled cut chain}
\label{proof:sec:representation}

Fix $n\geq2$ and $\beta>0$. Write $L_M$ for the polymer generator in
Theorem~\ref{thm:main}, and put $K(m)=\sum_jm_j$. We use a labelled
auxiliary process, but only through its stationary Dirichlet form.

Arrange $n$ labelled sites in cyclic order, with one bond from each site to
the next, and let $[n]$ also index these bonds. For $n=2$ these are two
distinct bonds between the two sites. Thus the construction concerns a
cyclic order and does not require a simple cycle graph. Let
\[
 \Omega_n=\{S\subseteq[n]:S\neq\varnothing\}.
\]
The elements of $S$ are the cut bonds. Define the generator
\begin{equation}\label{proof:eq:LB}
 (L_BF)(S)=\beta\sum_{i\notin S}\bigl(F(S\cup\{i\})-F(S)\bigr)
 +\frac{|S|-1}{n}\sum_{i\in S}\bigl(F(S\setminus\{i\})-F(S)\bigr).
\end{equation}
When $|S|=1$, the second sum is understood to be zero; the empty set is
never a state. The process is irreducible: one can pass from $S$ to $T$
by adding bonds until reaching $S\cup T$ and then deleting the bonds
outside $T$. Its reversible probability measure is
\begin{equation}\label{proof:eq:piB}
 \pi_B(S)=\frac1{Z_B}\frac{(\beta n)^{|S|}}{(|S|-1)!}.
\end{equation}
Indeed, if $|S|=k$ and $i\notin S$, then
$\pi_B(S\cup\{i\})/\pi_B(S)=\beta n/k$, whereas the reverse jump has
rate $k/n$.

Cutting the cyclic order at the bonds in $S$ produces $|S|$ segments.
Let $p(S)=(m_j)_{j\geq1}$ record their unordered lengths. In particular,
$\sum_jjm_j=n$ and $K(p(S))=|S|$.

\begin{proposition}\label{proof:prop:comparison}
The pushforward of $\pi_B$ under $p$ is
\begin{equation}\label{proof:eq:piM}
 \pi_M(m)=\frac1{Z_M}\frac{(\beta n)^{K(m)}}{\prod_jm_j!}.
\end{equation}
It is reversible for $L_M$. For every function $f$ on the polymer state
space, the lift $F=f\circ p$ satisfies
\begin{equation}\label{proof:eq:dirichletcomparison}
 \operatorname{Var}_{\pi_B}(F)=\operatorname{Var}_{\pi_M}(f),
 \qquad
 \mathcal E_B(F,F)=\mathcal E_M(f,f).
\end{equation}
Consequently $\gamma_M\geq\gamma_B$. More generally, any Poincar\'e
inequality for rotation-invariant functions on $\Omega_n$ holds, with
the same constant, for polymer observables.
\end{proposition}

\begin{proof}
If $m$ has $k$ segments, the number of cut sets mapping to $m$ is
\begin{equation}\label{proof:eq:multiplicity}
 N(m)=\frac{n(k-1)!}{\prod_jm_j!}.
\end{equation}
To prove this formula, count cut sets together with a distinguished
cut. On one hand there are $kN(m)$ such pairs. On the other hand the
distinguished cut has $n$ possible locations, and the ordered segment
lengths have $k!/\prod_jm_j!$ possible arrangements. This also proves
the formula when $k=1$ or $k=n$. Multiplying~\eqref{proof:eq:multiplicity}
by the weight per cut set in~\eqref{proof:eq:piB} gives
\eqref{proof:eq:piM}.

Conditional on $p(S)=m$, the cut set is uniform. Equivalently, after
distinguishing a uniformly chosen cut, its segment lengths occur in
uniform multiset order around the cycle. Adding an absent cut already
has the required fragmentation rates in every configuration: in a
segment of length $j$, the split into $i$ and $j-i$ has two possible
locations if $i<j/2$, and one if $i=j/2$. Thus its rate is
$2\beta m_j$ or $\beta m_j$, respectively.

For mergers, suppose first that $k\geq2$. At each of the $k$ boundaries,
the two neighbouring segment positions sample the multiset of lengths
without replacement. For $i\neq j$, the expected number of boundaries
between an $i$-segment and a $j$-segment is
\[
 k\left(\frac{m_i}{k}\frac{m_j}{k-1}
       +\frac{m_j}{k}\frac{m_i}{k-1}\right)
 =\frac{2m_im_j}{k-1}.
\]
For equal sizes it is $m_i(m_i-1)/(k-1)$. A boundary is removed at rate
$(k-1)/n$, so the averaged merger rates are $2m_im_j/n$ and
$m_i(m_i-1)/n$. These formulas include $k=2$: the same two segments
have two distinct boundaries, each with removal rate $1/n$.
At $k=1$ there is no merger.

These calculations prove the exact conditional identity
\begin{equation}\label{proof:eq:conditionalgenerator}
 \mathbb E_{\pi_B}\bigl[L_B(f\circ p)(S)\mid p(S)=m\bigr]
 =(L_Mf)(m).
\end{equation}
For two polymer functions $f,g$, reversibility of $L_B$ and
\eqref{proof:eq:conditionalgenerator} imply
\[
 \langle f,L_Mg\rangle_{\pi_M}
 =\langle f\circ p,L_B(g\circ p)\rangle_{\pi_B}
 =\langle L_Mf,g\rangle_{\pi_M}.
\]
Thus $L_M$ is reversible under $\pi_M$. Taking $g=f$, and using
$\mathcal E(f,f)=-\langle f,Lf\rangle$, proves the Dirichlet-form
identity. The variance identity follows from the pushforward measure.
Every polymer lift is invariant under a rotation of the labelled
cyclic order, so the final assertions follow from the variational
characterization of the spectral gap.
\end{proof}

\begin{remark}\label{proof:rem:averaging}
Equation~\eqref{proof:eq:conditionalgenerator} is a stationary
conditional average. The projection of the cut process need not be a
Markov chain when started from an individual cut configuration. The
proof uses only~\eqref{proof:eq:dirichletcomparison}; no assertion of
strong lumpability is needed.
\end{remark}

\section{Finite comparison and the cluster-count chain}
\label{proof:sec:count}

We first record a positive-vector comparison for finite killed
generators. Such comparisons are standard in the spectral analysis of
birth--death processes; see, for example, \cite{Zeifman1991}. We include
the argument to specify exactly what is used.

\begin{lemma}\label{proof:lem:positive}
Let $I$ be a finite set and let
\[
 (Ag)(x)=\sum_{y\neq x}r(x,y)(g(y)-g(x))-V(x)g(x),
 \qquad r(x,y)\geq0,\quad V(x)\geq0.
\]
Suppose the off-diagonal graph is connected. Assume positive weights
$\rho$ satisfy $\rho(x)r(x,y)=\rho(y)r(y,x)$ for all $x,y$.
If $\lambda_*$ is the smallest eigenvalue
of $-A$, then, for every strictly positive vector $h$,
\begin{equation}\label{proof:eq:positivecomparison}
 \min_x\frac{-(Ah)(x)}{h(x)}
 \ \leq\ \lambda_*\ \leq\
 \max_x\frac{-(Ah)(x)}{h(x)}.
\end{equation}
The conclusion also holds on a one-point state space.
\end{lemma}

\begin{proof}
The operator is self-adjoint in the inner product weighted by $\rho$,
and its quadratic form is
\[
 \langle g,-Ag\rangle_\rho
 =\frac12\sum_{x,y}\rho(x)r(x,y)(g(y)-g(x))^2
   +\sum_x\rho(x)V(x)g(x)^2.
\]
A unit vector minimizing this form can be replaced by its absolute
value without increasing the form. Hence a nonnegative minimizing
eigenvector $v$ exists. If $v(x)=0$, its eigenvalue equation implies
$\sum_y r(x,y)v(y)=0$, so all its neighbours also have value zero.
Connectedness shows that $v$ is strictly positive.

At a point $x$ where $v(x)/h(x)$ is minimal, the nonnegative
off-diagonal entries give
$(Av)(x)\geq(v(x)/h(x))(Ah)(x)$. Since $Av=-\lambda_*v$, this implies
$\lambda_*\leq-(Ah)(x)/h(x)$. A point where $v/h$ is maximal gives
the reverse inequality. This proves~\eqref{proof:eq:positivecomparison}.
For one state the assertion is immediate.
\end{proof}

The number of clusters is itself a Markov coordinate of the polymer
process. If it has value $k$, its total fragmentation and merger rates
are, respectively, $\beta(n-k)$ and $k(k-1)/n$. Write $Q$ for its
generator on $\{1,\ldots,n\}$:
\begin{equation}\label{proof:eq:Q}
 (Qg)(k)=\beta(n-k)(g(k+1)-g(k))
       +\frac{k(k-1)}n(g(k-1)-g(k)).
\end{equation}
The zero boundary coefficients suppress the nonexistent terms. The
stationary measure is the pushforward of $\pi_B$ under $S\mapsto|S|$.
Testing the polymer variational principle with count observables gives
\begin{equation}\label{proof:eq:countupper}
 \gamma_M\leq\gamma_Q.
\end{equation}

For $1\leq k\leq n-1$, define the discrete derivative
$(\partial g)(k)=g(k+1)-g(k)$. Direct expansion yields the intertwining
\begin{equation}\label{proof:eq:gradient}
 \partial Q=\mathcal A\partial,
\end{equation}
where $\mathcal A$ is the killed birth--death generator
\begin{align}
 (\mathcal Ah)(k)
 ={}&\beta(n-k-1)(h(k+1)-h(k))
     +\frac{k(k-1)}n(h(k-1)-h(k))\notag\\
 &-\left(\beta+\frac{2k}n\right)h(k).
 \label{proof:eq:derivativegenerator}
\end{align}
This finite identity is an instance of the birth--death
intertwinings studied in \cite{ChafaiJoulin2013}. Here it can also be
checked by comparing the three coefficients of $h(k-1),h(k),h(k+1)$.
The derivative is an isomorphism from functions modulo constants to
$\mathbb R^{n-1}$. Therefore the eigenvalues of $-\mathcal A$ are
exactly the nonzero eigenvalues of $-Q$, with multiplicity. In
particular its smallest eigenvalue is $\gamma_Q$.

Every finite birth--death matrix with positive interior rates is
reversible under positive weights, obtained recursively from detailed
balance. Lemma~\ref{proof:lem:positive} thus applies to $\mathcal A$;
for $n=2$ it is a one-point matrix. The choice $h=1$ gives
\begin{equation}\label{proof:eq:countbase}
 \gamma_Q\geq\beta+\frac2n.
\end{equation}

\begin{proposition}\label{proof:prop:upper}
For every $n\geq2$ and $\beta>0$,
\begin{equation}\label{proof:eq:uniformupper}
 \gamma_M\leq\gamma_Q
 \leq\beta+4\max\{\sqrt\beta,1/n\}
 \leq4\left(\beta+\sqrt\beta+\frac1n\right).
\end{equation}
\end{proposition}

\begin{proof}
Put $a=\max\{1,n\sqrt\beta\}$ and choose
$h(k)=((k+a)(k+a+1))^{-1}$. The adjacent ratios are
\[
 \frac{h(k+1)}{h(k)}=\frac{k+a}{k+a+2},\qquad
 \frac{h(k-1)}{h(k)}=\frac{k+a+1}{k+a-1}.
\]
Substitution into~\eqref{proof:eq:derivativegenerator} gives the exact
identity
\begin{equation}\label{proof:eq:uppercertificate}
 \frac{-(\mathcal Ah)(k)}{h(k)}
 =\frac{2\beta(n-k-1)}{k+a+2}+\beta
   +\frac{2ka}{n(k+a-1)}
 \leq\frac{2\beta n}{a}+\beta+\frac{2a}{n}.
\end{equation}
The last inequality uses $a\geq1$. If $a=n\sqrt\beta$, the last
expression is $\beta+4\sqrt\beta$. If $a=1$, then $\beta n\leq1/n$
and it is at most $\beta+4/n$. Apply the upper half of
Lemma~\ref{proof:lem:positive}, followed by
\eqref{proof:eq:countupper}.
\end{proof}

\section{Harmonic decomposition and rotation-invariant observables}
\label{proof:sec:harmonics}

The cut chain has the full permutation symmetry of $[n]$. Its
degree-one modes can be slow when $\beta$ is very small, but those
modes do not occur in polymer observables. We give the relevant
Boolean-lattice decomposition explicitly. Related harmonic
descriptions of functions on slices of the Boolean cube appear in
\cite{Filmus2016}.

For $0\leq j\leq n$, let $V_j$ be the space of real functions on the
$j$-subsets of $[n]$, with counting inner product. Set
$V_{-1}=V_{n+1}=\{0\}$. Define the upward and downward maps by
\[
 (Ua)(T)=\sum_{i\in T}a(T\setminus\{i\}),\qquad
 (Da)(R)=\sum_{i\notin R}a(R\cup\{i\}).
\]
Thus $U:V_j\to V_{j+1}$ and $D:V_j\to V_{j-1}$ are adjoints on
adjacent layers. On $V_j$, direct double counting gives
\begin{equation}\label{proof:eq:commutator}
 DU-UD=(n-2j)I.
\end{equation}
Indeed, the diagonal term of $DU$ counts $n-j$ extensions and that
of $UD$ counts $j$ deletions. Every other term exchanges one present
and one absent label and occurs once on each side.

For $0\leq\ell\leq\lfloor n/2\rfloor$, set
$H_\ell=\ker(D:V_\ell\to V_{\ell-1})$. If $a\in H_\ell$, define its
lift on a layer of $k$-subsets by
\begin{equation}\label{proof:eq:harmoniclift}
 F_a(S)=\sum_{\substack{T\subseteq S\\|T|=\ell}}a(T)
       =\frac{U^{k-\ell}a(S)}{(k-\ell)!}
       \quad(k\geq\ell),
\end{equation}
and put $F_a=0$ when $k<\ell$.

\begin{lemma}\label{proof:lem:layers}
For every $0\leq k\leq n$, the layer $V_k$ is the orthogonal direct
sum of the lifted spaces $H_\ell$, with
$0\leq\ell\leq\min\{k,n-k\}$. Their dimensions and lift norms are
\begin{equation}\label{proof:eq:dimensionnorm}
 \dim H_\ell=\binom n\ell-\binom n{\ell-1},
 \qquad
 \|F_a\|_{V_k}^2=\binom{n-2\ell}{k-\ell}\|a\|_{V_\ell}^2,
\end{equation}
where $\binom n{-1}=0$. The lift is injective for
$\ell\leq k\leq n-\ell$ and vanishes for $k>n-\ell$.
\end{lemma}

\begin{proof}
Taking an inner product in~\eqref{proof:eq:commutator} gives
\[
 \|Ua\|^2=\|Da\|^2+(n-2j)\|a\|^2\qquad(a\in V_j).
\]
Hence $U$ is injective for $j<n/2$, and its adjoint on the next layer
is surjective. This proves the dimension formula for $H_\ell$.

For $a\in H_\ell$, induction using the commutator gives
\begin{equation}\label{proof:eq:tower}
 DU^ja=j(n-2\ell-j+1)U^{j-1}a.
\end{equation}
For the induction step, commute $D$ past the last $U$ and use
\[
 (j-1)(n-2\ell-j+2)+n-2(\ell+j-1)
 =j(n-2\ell-j+1).
\]
Consequently
\[
 \|U^ja\|^2
 =j(n-2\ell-j+1)\|U^{j-1}a\|^2
 =\frac{j!(n-2\ell)!}{(n-2\ell-j)!}\|a\|^2
\]
for $0\leq j\leq n-2\ell$. Dividing by $(j!)^2$ proves the norm
formula and injectivity. At $j=n-2\ell+1$ the preceding norm
recursion gives zero, so all subsequent lifts vanish.

For orthogonality, let $a\in H_\ell$, $b\in H_h$, and $\ell<h$,
with both lifts supported on layer $k$. Move the $k-h$ upward maps
in the lift of $b$ across the inner product, and apply
\eqref{proof:eq:tower}. The resulting inner product is a scalar
multiple of $\langle U^{h-\ell}a,b\rangle$, which is zero because
$Db=0$. Finally the dimensions telescope:
\[
 \sum_{\ell=0}^{\min(k,n-k)}
 \left(\binom n\ell-\binom n{\ell-1}\right)
 =\binom n{\min(k,n-k)}=\binom nk.
\]
Thus the orthogonal subspaces exhaust $V_k$.
\end{proof}

On the nonempty Boolean lattice, the degree-$\ell$ sector consists
of all functions obtained by multiplying the lifts of a basis of
$H_\ell$ by arbitrary functions of the layer $k$. For degree zero
the layers are $1,\ldots,n$; for degree $\ell\geq1$ they are
$\ell,\ldots,n-\ell$. The layerwise decomposition in
Lemma~\ref{proof:lem:layers} is orthogonal also in $L^2(\pi_B)$,
because $\pi_B$ is constant on each layer.

\begin{lemma}\label{proof:lem:sectoraction}
Each degree sector is preserved by $L_B$. Degree zero is the count
generator $Q$. For $\ell\geq1$, each harmonic direction has radial
generator
\begin{align}
 (\mathcal L_\ell g)(k)
 ={}&\beta(n-k-\ell)(g(k+1)-g(k))\notag\\
 &+\frac{(k-1)(k-\ell)}n(g(k-1)-g(k))
  -\left(\ell\beta+\frac{\ell(k-1)}n\right)g(k),
 \label{proof:eq:sector}
\end{align}
on $\ell\leq k\leq n-\ell$. It is a reversible killed birth--death
generator. The spectrum of $-L_B$ is the union of the spectrum of
$-Q$ and the spectra of $-\mathcal L_\ell$, the latter repeated
$\dim H_\ell$ times.
\end{lemma}

\begin{proof}
The lift satisfies
\begin{equation}\label{proof:eq:neighbours}
 \sum_{i\in S}F_a(S\setminus\{i\})=(k-\ell)F_a(S),
 \qquad
 \sum_{i\notin S}F_a(S\cup\{i\})=(n-k-\ell)F_a(S).
\end{equation}
The first identity counts each $\ell$-subset of $S$ exactly $k-\ell$
times. For the second, the terms already contained in $S$ contribute
$(n-k)F_a(S)$. The remaining terms can be grouped as
\[
 \sum_{\substack{R\subseteq S\\|R|=\ell-1}}
 \sum_{i\notin S}a(R\cup\{i\}).
\]
By $Da=0$, the inner sum is the negative of the corresponding sum
over $i\in S\setminus R$. Summing over $R$ counts each
$\ell$-subset $\ell$ times, giving $-\ell F_a(S)$. For $\ell=0$
the identities hold directly.

Substitute~\eqref{proof:eq:neighbours} in~\eqref{proof:eq:LB} to obtain
\[
 L_B(F_ag)(S)=F_a(S)(\mathcal L_\ell g)(|S|).
\]
The same identities, with the zero lift outside its support, show
that there is no contribution outside this degree sector. Both
radial boundary coefficients in~\eqref{proof:eq:sector} are zero.

For an explicit self-adjointness check, the norm formula
\eqref{proof:eq:dimensionnorm} shows that the radial inner product,
up to a fixed positive multiplier, has weights
\begin{equation}\label{proof:eq:sectorweights}
 \nu_\ell(k)=\frac{(\beta n)^k}{(k-1)!}
               \binom{n-2\ell}{k-\ell}.
\end{equation}
They satisfy
\[
 \nu_\ell(k)\,\beta(n-k-\ell)
 =\nu_\ell(k+1)\,\frac{k(k+1-\ell)}n.
\]
The killing term is nonnegative. If the sector has more than one
layer, its interior birth and death rates are positive. Choosing an
orthogonal basis in each $H_\ell$ therefore yields orthogonal radial
copies of~\eqref{proof:eq:sector}. Completeness follows from
Lemma~\ref{proof:lem:layers}, proving the spectral assertion.
\end{proof}

\begin{proposition}\label{proof:prop:base}
Every rotation-invariant function $F$ on $\Omega_n$ satisfies
\begin{equation}\label{proof:eq:rotationpoincare}
 \mathcal E_B(F,F)\geq
 \left(\beta+\frac2n\right)\operatorname{Var}_{\pi_B}(F).
\end{equation}
In particular,
\begin{equation}\label{proof:eq:polymerbase}
 \gamma_M\geq\beta+\frac2n.
\end{equation}
\end{proposition}

\begin{proof}
For $\ell\geq2$, the killing in~\eqref{proof:eq:sector} is at least
\[
 \ell\beta+\frac{\ell(\ell-1)}n\geq2\beta+\frac2n.
\]
By Lemma~\ref{proof:lem:positive} with $h=1$, every eigenvalue of
$-\mathcal L_\ell$ is at least this value. In degree zero, every
nonconstant eigenvalue is at least $\beta+2/n$ by
\eqref{proof:eq:countbase}.

It remains to exclude degree one. The degree-one part on any layer
$1\leq k\leq n-1$ has the unique form
\[
 S\longmapsto\sum_{i\in S}a_i(k),\qquad\sum_{i=1}^na_i(k)=0.
\]
Uniqueness follows from the injectivity in
Lemma~\ref{proof:lem:layers}. All degree sectors and their
orthogonal complements are preserved by relabelling. Since the
inner product is also preserved, their orthogonal projections
commute with rotations. The degree-one projection of a
rotation-invariant $F$ is therefore rotation-invariant. Its unique
coefficient vector must be fixed by the cyclic permutation of all
$n$ labels, so its entries are all equal. Their sum is zero, and
the vector vanishes. The lift is already zero on layer $n$.

After subtracting its mean, $F$ consequently contains only the
nonconstant part of degree zero and degrees at least two. The
complete orthogonal spectral decomposition now proves
\eqref{proof:eq:rotationpoincare}. Apply
Proposition~\ref{proof:prop:comparison} to conclude
\eqref{proof:eq:polymerbase}.
\end{proof}

\section{A reciprocal certificate in the square-root regime}
\label{proof:sec:reciprocal}

\begin{proposition}\label{proof:prop:squareroot}
If $n\geq8$ and $4/n^2\leq\beta\leq1/16$, then
\begin{equation}\label{proof:eq:squareroot}
 \gamma_B\geq\frac{\sqrt\beta}{5},
 \qquad
 \gamma_M\geq\frac{\sqrt\beta}{5}.
\end{equation}
\end{proposition}

\begin{proof}
Put $a=n\sqrt\beta$, so that $2\leq a\leq n/4$. Use the positive
test vector $h(k)=1/(k+a)$, first for the count derivative
$\mathcal A$. Its exact quotient is
\begin{align}
 \frac{-(\mathcal Ah)(k)}{h(k)}
 &=\frac{\beta(n-k-1)}{k+a+1}
   -\frac{k(k-1)}{n(k+a-1)}+\beta+\frac{2k}n\notag\\
 &\geq\frac{\beta(n-k-1)}{k+a+1}+\beta+\frac kn.
 \label{proof:eq:countreciprocal}
\end{align}
If $k\leq a$, then $n-k-1\geq n/2$ and
$k+a+1\leq5a/2$. The birth term in the last line is therefore at
least $\beta n/(5a)=\sqrt\beta/5$. If $k\geq a$, its last term is
at least $a/n=\sqrt\beta$. Lemma~\ref{proof:lem:positive} proves
that every nonconstant degree-zero eigenvalue is at least
$\sqrt\beta/5$.

For degree $\ell\geq1$, the same test vector on
$\ell\leq k\leq n-\ell$ gives
\begin{equation}\label{proof:eq:sectorcertificate}
 \frac{-(\mathcal L_\ell h)(k)}{h(k)}
 =\frac{\beta(n-k-\ell)}{k+a+1}+\ell\beta
  +\frac{(k-1)((\ell-1)k+\ell a)}{n(k+a-1)}.
\end{equation}
For clarity, the last term combines the negative death contribution
$-(k-1)(k-\ell)/(n(k+a-1))$ with the killing
$\ell(k-1)/n$. Their numerator is
\[
 (k-1)\bigl(\ell(k+a-1)-(k-\ell)\bigr)
 =(k-1)((\ell-1)k+\ell a).
\]
All terms on the right of~\eqref{proof:eq:sectorcertificate} are
nonnegative. If $k\leq a$, then $\ell\leq k$ gives
$n-k-\ell\geq n-2a\geq n/2$, while $k+a+1\leq5a/2$.
The birth term is again at least $\sqrt\beta/5$. If $k\geq a$,
then $(\ell-1)k+\ell a\geq a$ and
\[
 \frac{k-1}{k+a-1}\geq\frac13,
 \qquad
 3(k-1)-(k+a-1)=2k-a-2\geq a-2\geq0.
\]
The last term in~\eqref{proof:eq:sectorcertificate} is at least
$a/(3n)=\sqrt\beta/3$. Another application of
Lemma~\ref{proof:lem:positive} bounds every eigenvalue of every
$-\mathcal L_\ell$ below by $\sqrt\beta/5$.

Lemma~\ref{proof:lem:sectoraction} exhausts the spectrum of the cut
chain, so $\gamma_B\geq\sqrt\beta/5$. The polymer inequality follows
from Proposition~\ref{proof:prop:comparison}.
\end{proof}

\section{Completion of the uniform bound}
\label{proof:sec:completion}

\begin{proof}[Proof of Theorem~\ref{thm:main}]
Write $s=\beta+\sqrt\beta+1/n$. Proposition~\ref{proof:prop:upper}
gives $\gamma_M\leq4s$. For the lower bound, use
\eqref{proof:eq:polymerbase} throughout and divide the parameters
into the following exhaustive cases.

If $\beta\geq1/16$, then $\sqrt\beta\leq4\beta$, whence
\[
 s\leq5\beta+\frac1n
 \leq5\left(\beta+\frac2n\right)\leq5\gamma_M.
\]
If $\beta\leq4/n^2$, then $\sqrt\beta\leq2/n$, so
\[
 s\leq\beta+\frac3n
 \leq\frac32\left(\beta+\frac2n\right)
 \leq\frac32\gamma_M.
\]
If $n<8$ and $\beta<1/16$, the same conclusion follows from
$\sqrt\beta<1/4\leq2/n$.

In every remaining case, $n\geq8$ and
$4/n^2<\beta<1/16$. Proposition~\ref{proof:prop:squareroot} gives
$\sqrt\beta\leq5\gamma_M$, while
\eqref{proof:eq:polymerbase} gives $\beta+1/n\leq\gamma_M$.
Adding these inequalities yields $s\leq6\gamma_M$.
The preceding weak inequalities cover all boundary values.
Therefore $s/6\leq\gamma_M\leq4s$ for every $n\geq2$ and
$\beta>0$. Both bounds are positive, so their reciprocals give the
stated relaxation-time bounds.
\end{proof}
